\documentclass[12pt]{article}
\usepackage[utf8]{inputenc}
\usepackage[T1]{fontenc}
\usepackage{hyperref}
\title{Adai Language Resources for K-12 Curriculum Development}
\author{Compiled by Claude Sonnet 5.5 (claude-sonnet-5-5) for Andr\'as Kornai}
\date{}
\begin{document}
\maketitle

\section{Introduction}
Adai (ISO 639-3 xad; Glottocode adai1235) is an extinct, unclassified language of northwestern Louisiana. It is known only from a list of about 275 words that John Sibley recorded in 1804. Some sources give about 250 words. Early scholars linked it to Caddoan, but that view is now doubted. Later work suggests the list is Adai with some Caddoan borrowings. The language may be an isolate. The Adai people were affected by Spanish colonization, disease and conflict, and surviving families near Nacogdoches merged into the Caddo. The Adai Caddo Indian Nation says the language was gone by the mid-1800s.

\section{Primary Educational Resources}
No primers, textbooks or school programmes for Adai were found.

\section{Linguistic and Reference Materials}
\begin{itemize}
\item Anthony P.\ Grant (1995), whose paper reproduces the whole Sibley vocabulary from the manuscript at the American Philosophical Society (Glottolog reference record): \url{https://glottolog.org/resource/reference/id/76121}
\item Wikipedia on the Adai language: \url{https://en.wikipedia.org/wiki/Adai_language}
\item Wikipedia on the Adai people: \url{https://en.wikipedia.org/wiki/Adai_people}
\item Numerals page with the Adai data (the page warns the text transcribed from Grant 1995 may be inaccurate because of poor printing): \url{https://lingweb.eva.mpg.de/channumerals/Adai.htm}
\item Caddo materials, American Philosophical Society Archives guide: \url{https://indigenousguide.amphilsoc.org/entry/40258}
\item University of Utah digital collection item with OCR text that came up in the search (content not verified): \url{https://collections.lib.utah.edu/details?id=356201}
\end{itemize}

\section{Religious and Community Resources}
The Adai Caddo Indian Nation has a web presence with a page on its people and history: \url{https://adaination.com/?p=1528}

\section{Digital and Audio Resources}
No audio recordings, corpora or software were found. The Sibley word list is the only linguistic data found.

\section{Curriculum Development Context}
The language has no speakers and survives in a single short word list of uncertain accuracy. A K-12 programme can only be a heritage and history unit, working with the Adai Caddo Indian Nation, using the Sibley vocabulary with caveats. Pronunciation is not recoverable from the sources found. Any use of the words should say so plainly.

\end{document}